% Justificación como entregable aparte (justificacion.tex).
% Sigue el material de clase de Taller de Investigación I (TecNM): impacto
% social, tecnológico, económico y ético, y viabilidad. Se omite el impacto
% ambiental porque el proyecto no tiene una relación directa con el ambiente
% ni con el uso de recursos naturales.
% La justificación no lleva introducción ni tablas (las retiradas están en
% tablas/). La viabilidad está en secciones/12-viabilidad.tex.
% Redacción en tercera persona.
% Los antecedentes (estudios previos) están en secciones/11-antecedentes.tex.
%
% En justificacion.tex los apartados son secciones; dentro de otro documento
% (problema-objetivos.tex) quedan como subsecciones de «Justificación».
\ifdefined\solojustificacion
\let\justsec\section
\let\justsubsec\subsection
\else
\section{Justificación}
\let\justsec\subsection
\let\justsubsec\subsubsection
\fi

\justsec{Impacto}

\justsubsec{Impacto Social}
La investigación puede beneficiar, en primer lugar, a los estudiantes, que
pierden tiempo en las filas de entrada y se exponen a llegar tarde a su
primera clase. Medir el tiempo de validación y la longitud de la fila antes
y durante el piloto permitirá saber si un cambio en el mecanismo de acceso
alivia esa situación o si la espera se debe a otras causas.

También puede beneficiar al personal de vigilancia, que revisa cientos de
credenciales en pocos minutos y cuya revisión depende solo de su atención en
los momentos de mayor afluencia. Por último, el Departamento de Control
Escolar y los directivos contarían con datos de ingreso por franja horaria
que hoy no existen, útiles para planear la apertura de accesos y el personal
necesario en cada horario. Como información complementaria, el registro de
intentos rechazados aportaría una primera medida de la presentación de
credenciales ajenas o vencidas, una preocupación de seguridad cuya magnitud
tampoco se conoce.

\justsubsec{Impacto Tecnológico}
Desde la dimensión tecnológica, el proyecto integra componentes de hardware
y software en un mismo sistema: un microcontrolador con cámara que captura y
decodifica el código, un controlador que libera el torniquete, un servidor
que valida el código y una base de datos con un panel web para consultar los
registros.

La evidencia disponible sobre el ESP32-CAM aborda la transmisión de video y
no la decodificación de códigos QR \parencite{nowroz2025esp32cam}. Las
pruebas del prototipo bajo distintas condiciones de iluminación, brillo de
pantalla y distancia permitirán conocer su tiempo de lectura y su tasa de
lecturas fallidas. Estos resultados, junto con la documentación del diseño,
podrán servir como precedente para otros proyectos del \tescha{} que
busquen automatizar procesos con componentes de bajo costo.

\justsubsec{Impacto Económico}
El sistema se apoya en el celular que ya posee el estudiante, por lo que no
es necesario emitir tarjetas ni credenciales adicionales, y el lector se
construye con componentes de bajo costo. Frente a otras tecnologías de
validación, como NFC o el reconocimiento facial, el código QR requiere menos
infraestructura. Sin embargo, cualquier ahorro en tiempo de trabajo del
personal o en costos de operación tendría que comprobarse en evaluaciones
posteriores; esta investigación solo aportará la información del piloto para
valorarlo.

\justsubsec{Impacto Ético}
El sistema trata datos personales de los estudiantes, como su identidad y
la hora de ingreso. Por ello, antes del piloto se identificarán los
requisitos de la legislación estatal \parencite{edomex2017datos} y general
\parencite{lgpdppso2025}, y se definirán con Control Escolar el aviso de
privacidad y el resguardo de los registros. A diferencia del reconocimiento
facial, el código QR no trata datos biométricos, y quien no pueda mostrar su
código seguirá el procedimiento manual actual.

% La viabilidad pasará al documento de metodología; solo el entregable
% justificacion.tex la conserva.
\ifdefined\solojustificacion
% Viabilidad del proyecto. Se retiró de la justificación de
% problema-objetivos.tex porque irá en el documento de metodología (aún no
% existe). Por ahora solo la carga justificacion.tex.
% Redacción en tercera persona.
\ifdefined\justsec\justsec{Viabilidad}\else\section{Viabilidad}\fi
El proyecto propone construir un prototipo de torniquete con un lector
ESP32-CAM que valide un código QR dinámico mostrado en el celular del
estudiante y registre cada ingreso en una base de datos consultable desde
un panel web. El prototipo se probará en un acceso peatonal piloto del
turno vespertino para evaluar su efecto en el tiempo de validación y en la
longitud de la fila.

La viabilidad técnica se sustenta en que los componentes del prototipo, el
ESP32-CAM, el controlador del torniquete, el servidor y la base de datos,
están disponibles comercialmente y corresponden a conocimientos que se han
adquirido en Ingeniería en Sistemas Computacionales. Ya se ha trabajado con
captura de imágenes en pruebas de banco durante la exploración de
alternativas, lo que reduce el riesgo de la etapa de desarrollo. Frente a
NFC, una comparación en transacciones móviles encontró que el QR se procesa
con mayor rapidez, aunque no incorpora cifrado propio
\parencite{harnaningrum2024comparison}; esa carencia se compensa con la firma
HMAC del código.

En cuanto a los recursos económicos, el bajo costo de los componentes
permite construir y probar el prototipo con recursos propios, sin depender
de una inversión institucional. En cuanto al tiempo, el trabajo se organiza
en cinco fases (diagnóstico, diseño de hardware y firmware, desarrollo del
servidor, piloto controlado y evaluación), y el estudio se limita a un solo
acceso peatonal y a la franja de mayor afluencia del turno vespertino, lo
que lo hace abarcable en el periodo escolar.

El acceso a la información y a los participantes es posible porque el
proyecto se desarrolla dentro del \tescha{}: se puede observar el acceso
piloto y aplicar la encuesta a estudiantes y las entrevistas al personal de
vigilancia y de Control Escolar. La instalación del prototipo en el acceso
requiere la autorización de las autoridades del plantel, que se solicitará
antes del piloto.

Queda por confirmar la viabilidad desde el lado del estudiante, es decir, si
cuenta con un celular propio, con conexión a internet y con batería
suficiente al llegar al plantel. Esta condición se indagará con la encuesta
del diagnóstico, y sus resultados indicarán si el sistema puede usarse de
forma general o si el procedimiento manual debe mantenerse para una parte de
los estudiantes.

\fi

\justsec{Conclusión}
Hoy no existen datos que midan el problema del acceso en el \tescha{}: se
desconoce cuánto tiempo espera un estudiante para entrar y qué longitud
alcanzan las filas. La investigación se justifica porque puede generar esa
información y contrastarla con el desempeño de un prototipo de bajo costo.

En conjunto, los resultados podrán aportar evidencia del propio plantel para
valorar si un sistema de acceso con códigos QR dinámicos ayuda a aliviar
las esperas del sistema actual y si conviene ampliar el piloto a otros
accesos. No obstante, esa posible ampliación y sus beneficios económicos
tendrían que confirmarse mediante evaluaciones posteriores.
